\section{Introduction}\label{sec:introduction}

Bernoulli bond percolation on an infinite graph can have two distinct
transitions. Infinite open clusters first appear at a critical parameter
$\pc$; uniqueness begins at a second threshold $\pu$. On a nonamenable
graph, the expanding geometry can keep infinite clusters apart even
after percolation has begun. Benjamini and Schramm conjectured that this
mechanism always produces a nonempty nonuniqueness phase when the graph
has finitely many vertex orbits
\cite[Conjecture~6]{BenjaminiSchramm1996}. We prove this conjecture for
Bernoulli bond percolation, together with the stronger strict inequality
for the $\ell^2$ operator threshold conjectured by Hutchcroft.
This is the equivalent $\ell^2$ formulation of his operator-threshold
conjecture \cite[Proposition~2.3 and Conjecture~6.1]{Hutchcroft2019}.

\subsection{The graph, the thresholds, and the main result}

Let $G=(V,E)$ be an infinite connected locally finite undirected graph.
It is \emph{quasi-transitive} if its automorphism group $\Aut(G)$ has
finitely many orbits on $V$. In particular its vertex degrees are
uniformly bounded. For a finite nonempty set $A\subset V$, write
\[
 \partial_V A=\{v\in V\setminus A:v\sim a\text{ for some }a\in A\},
 \qquad
 h_V(G)=\inf_{0<|A|<\infty}\frac{|\partial_V A|}{|A|}.
\]
The graph is \emph{nonamenable} if $h_V(G)>0$. For a bounded-degree
graph this is equivalent to positivity of the edge isoperimetric
constant. Loops and multiple edges may be allowed; a subdivision
reduction in \cref{pre:simple-reduction} permits the main proof to use
a simple graph.

Under $\Prob_p$, every unordered edge is independently open with
probability $p$. Write $C_x$ for the open cluster of $x$, and
$x\leftrightarrow y$ for connection by a finite open path. Paths of
length zero are allowed. If $N_\infty(p)$ is the number of infinite
open clusters, define
\begin{align*}
 \pc(G)&=\inf\{p\in[0,1]:\Prob_p(N_\infty(p)\ge1)>0\},\\
 \pu(G)&=\inf\{p\in[0,1]:\Prob_p(N_\infty(p)=1)=1\}.
\end{align*}
The two-point connection kernel is $T_p(x,y)=\Prob_p(x\leftrightarrow y)$.
We equip $V$ with counting measure and put
\[
 \ell^2(V)=\left\{f:V\to\mathbb C:\sum_{x\in V}|f(x)|^2<\infty\right\},
 \qquad (T_pf)(x)=\sum_{y\in V}T_p(x,y)f(y)
\]
initially for finitely supported $f$. Boundedness of $T_p$ means that
this action extends to a bounded operator on $\ell^2(V)$. Its threshold is
\[
 \ptwo(G)=\sup\{p\in[0,1]:T_p\text{ is bounded on }\ell^2(V)\}.
\]
All operator norms in this paper use counting measure, including
operators restricted to a vertex orbit. The weights introduced later
average the choice of root and define a finite trace; they do not
change the Hilbert-space measure.

\begin{theorem}\label{thm:main}
Let $G$ be an infinite connected locally finite quasi-transitive graph
with $h_V(G)>0$. For Bernoulli bond percolation on $G$,
\[
 \sup_{0\le p<\pc}\norm{T_p}_{2\to2}
       =\norm{T_{\pc}}_{2\to2}<\infty,
 \qquad \pc(G)<\ptwo(G)\le\pu(G).
\]
There exist $\pc<p_1<p_2<1$ such that, in the coupling obtained from
independent uniform edge labels, almost surely there are infinitely
many infinite open clusters simultaneously for every $p\in[p_1,p_2]$.
\end{theorem}

The theorem resolves the bond-percolation nonuniqueness conjecture of
Benjamini and Schramm positively, and proves Hutchcroft's stronger
operator-threshold conjecture in the same graph class. The bounds and
the interval may depend on $G$. The simultaneous conclusion is a
consequence of the strict gap, in accordance with the simultaneous
phase theorem of H\"aggstr\"om, Peres, and Schonmann
\cite[Theorem~1.3]{HaggstromPeresSchonmann1999}. We include a direct
proof for the compact interval asserted here.

\begin{corollary}[Arbitrary Cayley graphs]\label{cor:cayley}
Let $\Lambda$ be a nonamenable finitely generated group, and let $S$
be any finite symmetric generating set with the identity removed.
For Bernoulli bond percolation on the simple undirected Cayley graph
with edges $\{x,xs\}$, $x\in\Lambda$, $s\in S$, the critical connection
operator is bounded on counting-measure $\ell^2(\Lambda)$ and
\[
 \pc<\ptwo\le\pu.
\]
The simultaneous nonuniqueness interval in \cref{thm:main} exists for
every such fixed generating set. At every fixed $p\in(\pc,\pu)$,
there are almost surely infinitely many infinite open clusters.
\end{corollary}

\begin{proof}
Left translations act transitively on this locally finite graph.
The graph form of the F\o lner criterion gives positive isoperimetric
constant for every fixed finite generating set of a nonamenable group.
Apply \cref{thm:main}.
The final assertion follows from the phase theorem
\cite[Theorem~1.3]{HaggstromPeresSchonmann1999}.
\end{proof}

\subsection{Proof strategy}

For $p\in[0,1]$, write
\[
 \chi_x(p)=\E_p|C_x|=\sum_yT_p(x,y),\qquad
 \chi_{\max}(p)=\sup_x\chi_x(p),
\]
for the susceptibility and its supremum over vertices. These quantities
are finite for $p<\pc$ by sharpness (\cref{pre:sharpness}). The row sums
bound the operator norm, but that bound need not remain finite as
$p\uparrow\pc$. Our task is to control
$\norm{T_\pc}_{2\to2}$ independently of these row sums. The proof first
bounds their near-critical growth and then constructs a deformation
that reduces row sums while retaining a large operator norm.
The nonunimodular case follows from Hutchcroft's established operator
theorem. After the reductions in \cref{pre:section}, the new argument
therefore concerns a simple graph with a unimodular automorphism group.
Its critical clusters are finite by the critical-nonpercolation theorem
recalled there.

\paragraph{From local endpoint maps to critical moments.}
A local endpoint map is a randomized rule that either fails or selects
a vertex in the root's critical open cluster, with a fixed cap $M$ on
that whole cluster. The roots and selected vertices lie in one vertex
orbit. Let $u$ be its success probability, and let $K(x,y)$ be the
probability of selecting $y$ conditional on success at $x$.
If some critical moment of order
$0<\alpha<1/2$ were infinite, \cref{amp:section} would construct such
maps with $0<\norm K<1$ and
\[
 \frac{\log(1/u)+\log M}{\log(1/\norm K)}\longrightarrow0.
\]
Thus the endpoint kernel becomes dispersive much faster than the success
probability deteriorates or the cap grows.

The construction exchanges finite pieces between independent
percolation fields. The geometric input, proved in \cref{good:section},
holds outside a stretched-exponentially small event: after any vertex
deletion from a capped critical cluster, it bounds the number of bridges
separating two vertices in each remaining component that has an open
contact with the deleted set, in terms of its number of contacts.
An inverse comparison and a
joint second-moment bound in \cref{amp:section} control the multiplicity
of the exchanges and keep the output mass from concentrating on rare
fields. In \cref{mom:section}, each fixed map is run just below $\pc$.
Its small operator norm forces its iterates to send substantial endpoint
mass outside an independently sampled cluster whose law is weighted
by its size. Comparing the susceptibility in that deleted
domain with the full susceptibility gives a differential inequality.
A finite-exploration comparison with critical percolation then
contradicts the assumed divergent moment. This proves
\[
 \sup_x\E_\pc|C_x|^\alpha<\infty\quad(0<\alpha<1/2).
\]
A separate decision-tree argument in the same section converts these
moments into
\[
 \chi_{\max}(p)\le C_\eta(\pc-p)^{-1-\eta}
 \qquad\text{for every }\eta>0.
\]

\paragraph{A subpower triangle bound.}
The next quantity is the triangle diagram
\[
 \nabla_p(x)=\sum_{y,z}T_p(x,y)T_p(y,z)T_p(z,x).
\]
\Cref{trace:section} compares two explorations of a cluster, one in the
full graph and one after another cluster has been removed. The
difference of their centered edge scores---sums of tested edge bits
after subtracting their means---counts the closed contacts between
the clusters. Concentration of the scores and the fractional
moments above bound those contacts. A differential inequality for the
type average of the triangle diagram then yields
\[
 \sup_x\nabla_p(x)\le C_\eta(\pc-p)^{-\eta}
 \qquad\text{for every }\eta>0.
\]
The supremum and type-average forms are equivalent up to fixed constants
because there are finitely many vertex types. This subpower bound is the
remaining input to the deformation argument.

\paragraph{Corridor deformation and the spectral contradiction.}
A corridor is a simple open path whose internal vertices have no other
open incident edges. For fixed $p<\pc$, \cref{def:section} penalizes a
connection for each specified corridor all of whose edges are pivotal
for it. Raw paths of length $k$ from a symmetric lazy walk specify these
corridors by loop erasure; the penalty weights cancel the probabilities
of forcing their open edges and closed boundaries.

Nonamenability gives the walk an $\ell^2$ norm $\rho<1$. After deformation
time $t$, the norm of the connection kernel loses at most
$2t\rho^k\norm{T_p}^2$. Up to an explicit time-dependent factor, the
rate of decrease of its averaged row sum is bounded below by an average
of products of the deformed susceptibilities, or weighted row sums, at the two
corridor endpoints in the graph with the corridor interior deleted,
apart from an error
bounded by a constant times
$\chi_{\max}(p)^2\rho^k\sup_x\nabla_p(x)$.
\Cref{path:section} finds a positive proportion of raw paths for which
both endpoint row sums remain substantial. Geodesic segments
at the two ends provide escapes from the raw path, and random-walk
smoothing makes the effect of deleting the corridor small at the
escape endpoints.

If $T_\pc$ were unbounded, sprinkling would imply
$\norm{T_p}\ge c(\pc-p)^{-1}$. The preceding susceptibility bound
would make $\chi_{\max}(p)/\norm{T_p}$ subpower as well. In
\cref{con:section}, these estimates permit a length
$k=o(\log(1/(\pc-p)))$ for which the norm stays large while the
averaged row sum decreases quadratically in its own size. The row sum
would become too small to bound the surviving operator norm. This
contradiction proves critical boundedness. Sprinkling then gives
$\pc<\ptwo$, the uniqueness comparison gives $\ptwo\le\pu$, and a
finite-modification argument supplies the simultaneous interval in
\cref{thm:main}.

\subsection{Critical behavior and exponential connectivity}

The operator conclusion gives more than the separation of the two
percolation thresholds. It also determines the powers governing the
size and radius of critical clusters, and ensures that two-point
connections still decay exponentially in part of the supercritical
phase. We state these consequences using the following observables.
The triangle diagram $\nabla_p(v)$ and susceptibility $\chi_v(p)$
were defined above; either may be infinite. For $p\in[0,1]$ and
$v\in V$, write
\[
 \theta_v(p)=\Prob_p(|C_v|=\infty),
\]
and let $G[p]$ be the open spanning subgraph.
The extrinsic and intrinsic radii of the cluster of $v$ are
\[
 \operatorname{rad}_{\mathrm{ext}}(C_v)
    =\sup_{x\in C_v}\dist_G(v,x),\qquad
 \operatorname{rad}_{\mathrm{int}}(C_v)
    =\sup_{x\in C_v}\dist_{G[p]}(v,x),
\]
with values in $[0,\infty]$. Thus the intrinsic radius measures
distance along open paths, whereas the extrinsic radius uses the
ambient graph.

Say that connection probabilities at $p$ decay exponentially if
there are constants $0<A<\infty$ and $a>0$ such that
$T_p(x,y)\le A e^{-a\dist_G(x,y)}$ for every $x,y\in V$, and define
\[
 \pexp(G)=\sup\{p\in[0,1]:
       \text{connection probabilities at $p$ decay exponentially}\}.
\]
This is the exponential-connectivity threshold of
\cite[Section~2]{Hutchcroft2020L2}. We write $f\asymp g$ in a stated
limit when $c g\le f\le Cg$ sufficiently near that limit for some
constants $0<c\le C<\infty$.

\Needspace{25\baselineskip}
\begin{corollary}[Critical laws and exponential connectivity]
\label{cor:critical-laws}
Let $G$ be an infinite connected locally finite quasi-transitive graph
with $h_V(G)>0$. For Bernoulli bond percolation on $G$,
\begin{equation}\label{eq:critical-triangle}
 \nabla_{\pc}(v)\le\norm{T_{\pc}}_{2\to2}^{3}<\infty
 \qquad\text{for every }v\in V.
\end{equation}
For every $v\in V$,
\begin{align}
 \chi_v(p)&\asymp(\pc-p)^{-1} &&(p\uparrow\pc), \label{eq:meanfield-chi}\\
 \theta_v(p)&\asymp p-\pc &&(p\downarrow\pc), \label{eq:meanfield-theta}\\
 \Prob_{\pc}(|C_v|\ge n)&\asymp n^{-1/2} &&(n\to\infty), \label{eq:meanfield-size}\\
 \Prob_{\pc}\bigl(\operatorname{rad}_{\mathrm{int}}(C_v)\ge n\bigr)
    &\asymp n^{-1} &&(n\to\infty), \label{eq:meanfield-intrinsic}\\
 \Prob_{\pc}\bigl(\operatorname{rad}_{\mathrm{ext}}(C_v)\ge n\bigr)
    &\asymp n^{-1} &&(n\to\infty). \label{eq:meanfield-extrinsic}
\end{align}
Here $n$ tends through positive integers. The comparison constants
and the one-sided neighborhoods of $\pc$ may depend on $G$ and $v$.
Moreover,
\begin{equation}\label{eq:exponential-threshold}
 \pc(G)<\ptwo(G)\le\pexp(G).
\end{equation}
For each fixed $p\in[0,\ptwo(G))$, there are constants
$0<A_{G,p}<\infty$ and $a_{G,p}>0$ such that
\begin{equation}\label{eq:exponential-connectivity}
 T_p(x,y)\le A_{G,p}e^{-a_{G,p}\dist_G(x,y)}
 \qquad\text{for every }x,y\in V.
\end{equation}
\end{corollary}

The triangle-condition assertion and the strict inequality
$\pc<\pexp$ verify Conjectures~1.2 and~2.1 of
\cite{Hutchcroft2020L2}. In particular, exponential connectivity
persists at $\pc$ and throughout the nonempty interval
$(\pc,\ptwo)$, even though infinite clusters already exist there.
The power laws are applications of earlier conditional theorems:
the new input is that the operator condition holds on every graph in
the stated class. We give the deduction, including the graph
conventions and the distinction between intrinsic and extrinsic
radius, in \cref{cons:section}.

\subsection{Historical context and the methods used here}

The distinction between existence of an infinite cluster and its
uniqueness is central to the program initiated by Benjamini and Schramm
\cite{BenjaminiSchramm1996}. Their formulation uses almost transitive
graphs, meaning the quasi-transitive graphs considered here.
Their standing model is site percolation, and they explicitly extend
their nonplanar questions to bond percolation on page~72. A regular
tree of degree at least three supplies the basic example of separated
thresholds. The difficulty is to extract an equally robust separation
from nonamenability alone, without a tree structure or an additional
geometric hypothesis.

Several earlier results identify additional geometry that guarantees
nonuniqueness. Benjamini and Schramm proved $0<\pc<\pu<1$ for
one-ended nonamenable transitive planar graphs, for both bond and site
percolation \cite[Theorem~1.1]{BenjaminiSchramm2001}. Their result
followed Lalley's work on site percolation on canonical graphs of
cocompact Fuchsian groups \cite{Lalley1998}. Schonmann established
multiple phase transitions in a quantitatively highly nonamenable
regime \cite{Schonmann2001}. Nachmias and Peres proved nonuniqueness
and triangle estimates for nonamenable transitive graphs of sufficiently large
girth relative to their spectral radius \cite{NachmiasPeres2012}.
These results explain the usefulness of geometric separation and
random-walk decay, while imposing conditions beyond nonamenability.

A different approach uses harmonic functions of finite Dirichlet energy.
Gaboriau proved $\pc<\pu$ for unimodular transitive graphs admitting a
nonconstant harmonic Dirichlet function, using $\ell^2$ Betti numbers,
and extended the argument to quasi-transitive graphs
\cite[Theorem~0.6 and Section~5]{Gaboriau2005}. This supplies a further
class in which nonuniqueness follows from a hypothesis stronger than
nonamenability alone.

For Cayley graphs, the quantifier over generating sets is significant.
Pak and Smirnova-Nagnibeda proved nonuniqueness for a suitable generating
multiset \cite{PakSmirnovaNagnibeda2000}. Thom proved that a suitable ordinary
finite symmetric generating set exists \cite[Corollary~8]{Thom2015}. These results
permit a choice of generators; \cref{cor:cayley} treats an arbitrary
prescribed generating set. The quasi-transitive theorem also permits
several vertex orbits and nontrivial vertex stabilizers, for which group
coordinates are not available.

Critical nonpercolation is a separate established input.
Benjamini, Lyons, Peres, and Schramm proved that critical percolation on
every nonamenable Cayley graph has no infinite clusters, and explained
the extension to unimodular quasi-transitive graphs
\cite[Theorem~1.1 and Section~2]{BLPS1999}. Hutchcroft established this conclusion on
every connected locally finite quasi-transitive graph of exponential
growth \cite[Theorem~1.2]{Hutchcroft2016}. Positive isoperimetry implies
exponential growth, so the latter theorem gives the finite critical
clusters on which our local constructions operate.

The behavior of infinite clusters in a common coupling has its own
history. H\"aggstr\"om and Peres proved uniqueness monotonicity and
simultaneous phase statements under unimodularity
\cite{HaggstromPeres1999}. Schonmann established stability and
uniqueness monotonicity for all quasi-transitive graphs
\cite{Schonmann1999}. H\"aggstr\"om, Peres, and Schonmann then proved
simultaneous phase counts away from the two threshold values in that
generality \cite[Theorem~1.3]{HaggstromPeresSchonmann1999}. These
results identify the probabilistic consequences of a strict threshold
gap; establishing that gap is the task here.

Hutchcroft introduced and developed the operator-threshold framework
used here, proved $\pc<\ptwo$ for nonamenable Gromov hyperbolic
quasi-transitive graphs, and proved it for graphs admitting a
quasi-transitive nonunimodular automorphism group
\cite[Section~2 and Theorems~2.1 and~2.9]{Hutchcroft2019}.
The nonunimodular result builds on his tilted approach to
nonuniqueness and mean-field criticality
\cite[Theorem~1.2]{Hutchcroft2020Nonunimodular}; the stronger operator
form used here is Theorem~2.9 of \cite{Hutchcroft2019}.
Choi and Seo subsequently obtained nonuniqueness and the triangle
condition for every Cayley graph of a finitely generated acylindrically
hyperbolic group by extending Hutchcroft's geometric criteria
\cite[Theorems~A and~B]{ChoiSeo2025}.
The nonunimodular theorem is used directly in our reduction. The new
argument therefore concerns a unimodular action. Its averaging over
vertex types uses the inverse stabilizer-volume weights from the
mass-transport theory of Benjamini, Lyons, Peres, and Schramm
\cite[Corollary~3.5]{BLPS1999Invariant}. The corresponding root law
and operator trace fit the general framework of Aldous and Lyons
\cite[Theorem~3.1 and Section~5]{AldousLyons2007}. We prove the
specific transport and trace identities needed below.

Hutchcroft developed the critical and supercritical consequences of
the operator condition in \cite{Hutchcroft2020L2}. His
Theorem~2.2 compares the operator and exponential-connectivity
thresholds, and his Theorem~3.1 gives the critical extrinsic-radius
law under the strict operator gap. The critical intrinsic-radius
upper bound has its origin in the lattice theorem of Kozma and
Nachmias \cite[Theorem~1.2(ii)]{KozmaNachmias2009}.
Hutchcroft's later finite-cluster theorem supplies both radius
estimates in the quasi-transitive setting under the operator
condition \cite[Theorem~1.2]{Hutchcroft2022FiniteClusters}.
These are the conditional geometric results applied in
\cref{cons:section}.
The upper comparison $\ptwo\le\pu$ can be an equality.
Hutchcroft and Pan proved $\ptwo=\pu$ for the Cartesian
product of a regular tree of degree at least three with an infinite
amenable Cayley graph
\cite[Theorem~1.2]{HutchcroftPan2024Dimension}. Thus the strict lower
comparison established here is compatible with a sharp upper endpoint.

Our probabilistic estimates use Harris's association inequality and
the van den Berg--Kesten inequality for disjoint open witnesses
\cite{Harris1960,VanDenBergKesten1985}. Subcritical susceptibility
finiteness is obtained by adapting the finite-set sharpness argument
of Duminil-Copin and Tassion \cite[Section~1]{DuminilCopinTassion2016}
to finitely many vertex types. Sharpness originates in the independent
work of Aizenman--Barsky and Menshikov
\cite{AizenmanBarsky1987,Menshikov1986}; Antunovi\'c and Veseli\'c
proved the full quasi-transitive susceptibility and exponential-tail
statements \cite{AntunovicVeselic2008}. The subsequent conversion of critical
cluster-size estimates into near-critical susceptibility estimates
uses independent vertex marks and a decision-tree covariance
inequality of O'Donnell, Saks, Schramm, and Servedio
\cite[Section~3.3, Equation~(6)]{ODonnellSaksSchrammServedio2005}.
This is the same critical-tail-to-susceptibility principle developed
by Hutchcroft \cite[Theorem~1.1]{Hutchcroft2020Exponents}; we give
the needed proof uniformly over vertex types and induced domains.
The sprinkling comparison that finally opens an interval above
$\pc$ is the connection-operator expansion of
\cite[Lemma~2.4]{Hutchcroft2019}, also proved here.

The contact estimates also have a close methodological antecedent in
Hutchcroft's two-ghost inequalities and their critical-tail refinements
\cite{Hutchcroft2020Locality,Hutchcroft2021PowerLaw}. Those arguments
use independent marks, cluster exploration, and centered edge scores,
following the fluctuation method of Aizenman, Kesten, and Newman
\cite{AizenmanKestenNewman1987}. Here we compare a full-cluster
exploration with an exploration in a deleted domain to control the
number of closed contacts between two clusters. The triangle diagram
and positive connection matrix belong to the tree-graph framework of
Aizenman and Newman \cite{AizenmanNewman1984}; the triangle-condition
work of Barsky--Aizenman and the subsequent quantitative development
by Hutchcroft provide further context
\cite{BarskyAizenman1991,Hutchcroft2022Triangle}. Our required subpower
trace bound and its use in the corridor comparison are proved below.

Finite modifications that join separated clusters, create pivotal edges,
and control the multiplicity of their preimages already appear in Kozma's
work on products of trees \cite[Section~1.4 and Lemmas~7--11]{Kozma2011}.
Our local endpoint maps exchange finite pieces between independent
percolation fields; the deletion-uniform bridge bounds and paired inverse
estimate needed for these exchanges are proved here. Together with the
deformation that penalizes forced pivotal corridors, these maps let us
compare the spreading of a connection kernel in $\ell^2$ with the total
mass of its rows. Both constructions are formulated through covariant
kernels on vertex orbits.
